% GENERATED FILE. Edit canonical lessons, then run npm run generate:books.
% canonical-source-hash: fnv1a64:1d0fa200ac48dd8f
% canonical-lessons: TA-C503-toli, TA-C503-anniyar, TA-C503-arimukam, TA-C503-upacari, TA-C503-piriyavitai

\chapter{Friends and Strangers}
\label{ch:ta-friends-and-strangers}
\hlchaptermodality{503}

\begin{chapteropening}
\textbf{By the end of this chapter:} \emph{I can talk about friends, strangers, introductions, hosting guests and farewells.}

Complete the last lesson of chapter 503: I can talk about friends, strangers, introductions, hosting guests and farewells.
\end{chapteropening}

\section[\texorpdfstring{\ta{தோழி} (tōḻi)}{தோழி (tōḻi)}]{\ta{தோழி} (tōḻi) — a female friend}

\label{lesson:TA-C503-toli}



\begin{quote}
Before the new one: say the Tamil for to stroke, then the Tamil for a bond.
\end{quote}

\subsection*{You'll want to know: \ta{தோழி}}
\textbf{\ta{தோழி}} — \emph{tōḻi} — \textquotedblleft{}a female friend\textquotedblright{}.

\begin{quote}
\ta{என்} \ta{தோழியும்} \ta{நானும்} \ta{ஒரே} \ta{பள்ளியில்} \ta{படித்தோம்}.
\end{quote}

\begin{quote}
\emph{eṉ tōḻiyum nāṉum orē paḷḷiyil paṭittōm.}
\end{quote}

\begin{quote}
My friend and I studied at the same school.
\end{quote}

A woman’s or girl’s friend; the man’s word is \ta{தோழன்}.

\begin{cousinweb}
\textbf{\ta{தோழி}} (\emph{tōḻi}) is native Dravidian (DEDR 3563), beside \ta{தோழன்} \emph{tōḻaṉ}, \textquotedblleft{}friend\textquotedblright{}.
\end{cousinweb}

\begin{grammarlens}[title={What you've built}]
One more word for describing what you lived and felt.
\end{grammarlens}

\subsection*{Your turn}
\begin{itemize}
\raggedright
  \item \emph{Say it:} \emph{tōḻi}
  \item \emph{Say it:} \emph{tōḻi}, once more
  \item \emph{Say it:} \emph{eṉ tōḻiyum nāṉum orē paḷḷiyil paṭittōm.} — My friend and I studied at the same school.
  \item \emph{Recall:} say the Tamil for to stroke, then the Tamil for a bond, then say \emph{tōḻi} again
\end{itemize}

\begin{tcolorbox}[breakable,colback=teal!4,colframe=teal!35!black,title={Before you move on}]
What does \ta{தோழி} mean? (\textquotedblleft{}A female friend\textquotedblright{}.) Say it once more.
\end{tcolorbox}
\section[\texorpdfstring{\ta{அந்நியர்} (anniyar)}{அந்நியர் (anniyar)}]{\ta{அந்நியர்} (anniyar) — a stranger; an outsider}

\label{lesson:TA-C503-anniyar}



\begin{quote}
Before the new one: say the Tamil for a bond, then the Tamil for a female friend.
\end{quote}

\subsection*{You'll want to know: \ta{அந்நியர்}}
\textbf{\ta{அந்நியர்}} — \emph{anniyar} — \textquotedblleft{}a stranger; an outsider\textquotedblright{}.

\begin{quote}
\ta{அந்நியர்களிடம்} \ta{குழந்தைகள்} \ta{பேசக்} \ta{கூடாது}.
\end{quote}

\begin{quote}
\emph{anniyarkaḷiṭam kuḻantaikaḷ pēcak kūṭātu.}
\end{quote}

\begin{quote}
Children must not talk to strangers.
\end{quote}

Plural: \ta{அந்நியர்கள்}, \emph{anniyarkaḷ}.

Respectful -\ta{அர்}; plural \ta{அந்நியர்கள்}.

\begin{cousinweb}
\textbf{\ta{அந்நியர்}} (\emph{anniyar}) is borrowed from Sanskrit \emph{anya}, \textquotedblleft{}other\textquotedblright{} (CDIAL 399), with the Tamil -\ta{அர்} of a person.
\end{cousinweb}

\begin{grammarlens}[title={What you've built}]
One more word for describing what you lived and felt.
\end{grammarlens}

\subsection*{Your turn}
\begin{itemize}
\raggedright
  \item \emph{Say it:} \emph{anniyar}
  \item \emph{Say it:} \emph{anniyar}, once more
  \item \emph{Say it:} \emph{anniyarkaḷiṭam kuḻantaikaḷ pēcak kūṭātu.} — Children must not talk to strangers.
  \item \emph{Recall:} say the Tamil for a bond, then the Tamil for a female friend, then say \emph{anniyar} again
\end{itemize}

\begin{tcolorbox}[breakable,colback=teal!4,colframe=teal!35!black,title={Before you move on}]
What does \ta{அந்நியர்} mean? (\textquotedblleft{}A stranger; an outsider\textquotedblright{}.) Say it once more.
\end{tcolorbox}
\section[\texorpdfstring{\ta{அறிமுகம்} (aṟimukam)}{அறிமுகம் (aṟimukam)}]{\ta{அறிமுகம்} (aṟimukam) — an introduction; an acquaintance}

\label{lesson:TA-C503-arimukam}



\begin{quote}
Before the new one: say the Tamil for a female friend, then the Tamil for a stranger.
\end{quote}

\subsection*{You'll want to know: \ta{அறிமுகம்}}
\textbf{\ta{அறிமுகம்}} — \emph{aṟimukam} — \textquotedblleft{}an introduction; an acquaintance\textquotedblright{}.

\begin{quote}
\ta{ஒரு} \ta{நண்பர்} \ta{மூலம்} \ta{எனக்கு} \ta{அவருடைய} \ta{அறிமுகம்} \ta{கிடைத்தது}.
\end{quote}

\begin{quote}
\emph{oru naṇpar mūlam eṉakku avaruṭaiya aṟimukam kiṭaittatu.}
\end{quote}

\begin{quote}
I got to know him through a friend.
\end{quote}

\ta{அறிமுகம்} \ta{செய்}, to introduce.

\begin{cousinweb}
\textbf{\ta{அறிமுகம்}} (\emph{aṟimukam}) is a native compound of \ta{அறி} \emph{aṟi}, \textquotedblleft{}to know\textquotedblright{} (DEDR 314), and \ta{முகம்} \emph{mukam}, \textquotedblleft{}face\textquotedblright{} (DEDR 4889).
\end{cousinweb}

\begin{grammarlens}[title={What you've built}]
One more word for describing what you lived and felt.
\end{grammarlens}

\subsection*{Your turn}
\begin{itemize}
\raggedright
  \item \emph{Say it:} \emph{aṟimukam}
  \item \emph{Say it:} \emph{aṟimukam}, once more
  \item \emph{Say it:} \emph{oru naṇpar mūlam eṉakku avaruṭaiya aṟimukam kiṭaittatu.} — I got to know him through a friend.
  \item \emph{Recall:} say the Tamil for a female friend, then the Tamil for a stranger, then say \emph{aṟimukam} again
\end{itemize}

\begin{tcolorbox}[breakable,colback=teal!4,colframe=teal!35!black,title={Before you move on}]
What does \ta{அறிமுகம்} mean? (\textquotedblleft{}An introduction; an acquaintance\textquotedblright{}.) Say it once more.
\end{tcolorbox}
\section[\texorpdfstring{\ta{உபசரி} (upacari)}{உபசரி (upacari)}]{\ta{உபசரி} (upacari) — to host; to look after a guest}

\label{lesson:TA-C503-upacari}



\begin{quote}
Before the new one: say the Tamil for a stranger, then the Tamil for an introduction.
\end{quote}

\subsection*{You'll want to know: \ta{உபசரி}}
\textbf{\ta{உபசரி}} — \emph{upacari} — \textquotedblleft{}to host; to look after a guest\textquotedblright{}.

\begin{quote}
\ta{வீட்டுக்கு} \ta{வந்த} \ta{விருந்தினர்களை} \ta{அம்மா} \ta{நன்றாக} \ta{உபசரித்தாள்}.
\end{quote}

\begin{quote}
\emph{vīṭṭukku vanta viruntiṉarkaḷai ammā naṉṟāka upacarittāḷ.}
\end{quote}

\begin{quote}
Mother looked after the guests who came home very well.
\end{quote}

Past (I …): \ta{உபசரித்தேன்}, \emph{upacarittēṉ}.

The guest takes -\ta{ஐ}.

\begin{cousinweb}
\textbf{\ta{உபசரி}} (\emph{upacari}) is borrowed from Sanskrit \emph{upacāra}, \textquotedblleft{}approach, behaviour, service\textquotedblright{} (CDIAL 2163), made into a Tamil verb in -\ta{இ}.
\end{cousinweb}

\begin{grammarlens}[title={What you've built}]
One more word for describing what you lived and felt.
\end{grammarlens}

\subsection*{Your turn}
\begin{itemize}
\raggedright
  \item \emph{Say it:} \emph{upacari}
  \item \emph{Say it:} \emph{upacari}, once more
  \item \emph{Say it:} \emph{vīṭṭukku vanta viruntiṉarkaḷai ammā naṉṟāka upacarittāḷ.} — Mother looked after the guests who came home very well.
  \item \emph{Recall:} say the Tamil for a stranger, then the Tamil for an introduction, then say \emph{upacari} again
\end{itemize}

\begin{tcolorbox}[breakable,colback=teal!4,colframe=teal!35!black,title={Before you move on}]
What does \ta{உபசரி} mean? (\textquotedblleft{}To host; to look after a guest\textquotedblright{}.) Say it once more.
\end{tcolorbox}
\section[\texorpdfstring{\ta{பிரியாவிடை} (piriyāviṭai)}{பிரியாவிடை (piriyāviṭai)}]{\ta{பிரியாவிடை} (piriyāviṭai) — a farewell}

\label{lesson:TA-C503-piriyavitai}



\begin{quote}
Before the new one: say the Tamil for an introduction, then the Tamil for to host.
\end{quote}

\subsection*{You'll want to know: \ta{பிரியாவிடை}}
\textbf{\ta{பிரியாவிடை}} — \emph{piriyāviṭai} — \textquotedblleft{}a farewell\textquotedblright{}.

\begin{quote}
\ta{வேலையை} \ta{விட்டுப்} \ta{போகும்} \ta{ஆசிரியருக்கு} \ta{மாணவர்கள்} \ta{பிரியாவிடை} \ta{கொடுத்தார்கள்}.
\end{quote}

\begin{quote}
\emph{vēlaiyai viṭṭup pōkum āciriyarukku māṇavarkaḷ piriyāviṭai koṭuttārkaḷ.}
\end{quote}

\begin{quote}
The students gave a farewell to the teacher who was leaving.
\end{quote}

\ta{பிரியாவிடை} \ta{கொடு}, to give a send-off.

\begin{cousinweb}
\textbf{\ta{பிரியாவிடை}} (\emph{piriyāviṭai}) is a native compound of the negative of \ta{பிரி} \emph{piri}, \textquotedblleft{}to part\textquotedblright{} (DEDR 4176), and \ta{விடை}, \textquotedblleft{}leave, a send-off\textquotedblright{}.
\end{cousinweb}

\begin{grammarlens}[title={What you've built}]
One more word for describing what you lived and felt.
\end{grammarlens}

\subsection*{Your turn}
\begin{itemize}
\raggedright
  \item \emph{Say it:} \emph{piriyāviṭai}
  \item \emph{Say it:} \emph{piriyāviṭai}, once more
  \item \emph{Say it:} \emph{vēlaiyai viṭṭup pōkum āciriyarukku māṇavarkaḷ piriyāviṭai koṭuttārkaḷ.} — The students gave a farewell to the teacher who was leaving.
  \item \emph{Recall:} say the Tamil for an introduction, then the Tamil for to host, then say \emph{piriyāviṭai} again
\end{itemize}

\begin{tcolorbox}[breakable,colback=teal!4,colframe=teal!35!black,title={Before you move on}]
What does \ta{பிரியாவிடை} mean? (\textquotedblleft{}A farewell\textquotedblright{}.) Say it once more.
\end{tcolorbox}
